\documentclass[12pt]{article}

% Packages
\usepackage{times}
\usepackage[utf8]{inputenc}
\usepackage{geometry}
\geometry{a4paper, margin=1in}
\usepackage{hyperref}
\usepackage{enumitem}

% Title Page Information
\title{\textbf{Software Requirements Specification} \\ for \\ \textbf{NexusNotes}}
\author{
    \textbf{Project Team: Code Club} \\
    Ramasani Adi Pranav - 24CSB0A57 \\
    Syed Mohammad Sami - 24CSB0A77 \\
    Janapala Harsha Vardhan Reddy - 24CSB0A28
}
\date{February 20, 2026}

\begin{document}

\maketitle
\thispagestyle{empty}
\newpage

\tableofcontents
\newpage

\section{Introduction}

\subsection{Purpose}
The purpose of this document is to define the high-level objectives, scope, and strategic direction of the Nexus Notes platform.
It provides a comprehensive overview of the system's goals, target users, stakeholders, and core functionalities.
This document serves as a foundational reference for developers, academic stakeholders, and institutional authorities by clearly outlining the problem domain and the proposed solution.
It also ensures a shared understanding of the system vision among all participants involved in the project lifecycle.

\subsection{Document Conventions}
This document was created based on the IEEE template for System Requirement Specification Documents.
The requirements are prioritized as High, Medium, or Low based on their critical impact on the system's core functionalities.

\subsection{Intended Audience}

\begin{itemize}
    \item \textbf{Students:} Access, contribute, and utilize academic resources.
    \item \textbf{Faculty Members:} Verify and validate academic content.
    \item \textbf{Nexus Notes Administrators:} Manage platform operations, user roles, and content governance.
    \item \textbf{Developers:} Design, develop, and maintain the system.
    \item \textbf{Academic Coordinators:} Align content with curriculum requirements.
\end{itemize}

\subsection{Product Scope}
The fragmentation of academic resources across multiple platforms creates significant inefficiencies for students seeking reliable study materials, often resulting in redundant efforts and information silos.
Nexus Notes is a web-based centralized academic resource management system designed for educational institutions to resolve this issue.
The platform aims to unify fragmented academic resources into a structured digital repository organized subject-wise and module-wise.
The proposed platform assigns dedicated digital repositories to every curriculum subject, segmented into specialized subsections for lecture notes, critical question banks, capstone project archives, and emerging technology trends.
The system integrates dynamic collaborative features, including a contribution mechanism for peer-to-peer resource sharing, a faculty verification layer to ensure content accuracy, and a categorization logic that prioritizes high-quality content based on user engagement.
Furthermore, to bridge the gap between theoretical coursework and industry demands, the system incorporates a real-time Tech-Trend feed that aggregates relevant technological advancements for each specific domain.
The scope of the current system does not include full-fledged learning management functionalities such as automated grading systems, examination modules, or institutional administrative processes.

\subsection{References}
\begin{itemize}
    \item Vision Document, Nexus Notes, Centralized Academic Resource Platform, Version 1.0, January 30, 2026.
    \item Use Cases for Nexus Notes, Version 1.0, February 13, 2026.
\end{itemize}

\section{Overall Description}

\subsection{Product Perspective}
Nexus Notes is designed as a centralized academic knowledge platform that integrates students, faculty, and academic resources into a unified digital ecosystem.
The system acts as a bridge between fragmented academic materials and structured learning processes.
It complements existing institutional systems by providing an organized repository for academic content, collaboration features, and real-time technology updates.
Nexus Notes operates as a web-based platform accessible to all authorized users within an academic institution.

\subsection{Product Functions}
The core capabilities of Nexus Notes include:
\begin{itemize}
    \item \textbf{Centralized Repository:} Provides subject-wise digital repositories for academic resources.
    \item \textbf{Collaborative Sharing:} Enables students and faculty to contribute and share content.
    \item \textbf{Content Verification:} Allows faculty members to validate and approve academic materials.
    \item \textbf{Intelligent Categorization:} Organizes content based on relevance and usage.
    \item \textbf{Advanced Search:} Supports fast and accurate retrieval of academic resources.
    \item \textbf{Tech-Trend Integration:} Displays real-time technology trends related to academic domains.
\end{itemize}

\subsection{User Classes and Characteristics}

\begin{itemize}
    \item \textbf{Primary Users:} Students who access and share academic resources. Students require easy access to reliable and structured academic resources to enhance learning efficiency.

    \item \textbf{Secondary Users:} Faculty members who review and verify content. Faculty members focus on ensuring content accuracy and academic relevance.

    \item \textbf{Administrative Users:} Nexus Notes administrators responsible for managing the platform, moderating content, and controlling user permissions.

    \item \textbf{Contributors:} Users who upload and curate academic content.
\end{itemize}

\subsection{Operating Environment}
Users interact with Nexus Notes through web browsers on devices such as laptops, tablets, and smartphones.
The platform is accessed in various environments including classrooms, libraries, hostels, homes, and laboratories.
The system supports collaborative academic activities with continuous updates of digital content.

\subsection{Design and Implementation Constraints}
The Nexus Notes platform is subject to technical limitations such as server capacity, storage availability, network bandwidth, and system scalability.
The performance of the system depends on reliable internet connectivity and compatible hardware devices.
The system must ensure data security, user privacy, and protection of academic content.
Compliance with institutional data policies and ethical standards is required.

\subsection{User Documentation}
\begin{itemize}
    \item \textbf{User Manual:} A user manual is provided to guide students, faculty, and administrators in using the Nexus Notes platform.
    \item \textbf{Administrator Guide:} An administrator guide describes system configuration, user management, and content moderation procedures.
    \item \textbf{Developer Documentation:} Technical documentation supports system maintenance, future enhancements, and integration with other institutional systems.
\end{itemize}

\subsection{Assumptions and Dependencies}
It is assumed that users have access to stable internet connectivity and compatible digital devices.
The platform depends on institutional support for adoption and faculty participation for content verification.
Additionally, the effectiveness of the system relies on active user engagement and continuous content contribution.

\section{External Interface Requirements}

\subsection{User Interfaces}
Users interact with Nexus Notes through modern web browsers.
The interface is divided into role-based dashboards:
\begin{itemize}
    \item \textbf{Student Dashboard}
    \item \textbf{Administrator Console}
\end{itemize}

\subsection{Hardware Interfaces}
The platform supports cross-platform access via modern web browsers.

\subsection{Software Interfaces}
The platform requires a web server, database management system, and stable internet connectivity.
It may integrate with institutional authentication systems and external tech APIs.

\subsection{Communications Interfaces}
Nexus Notes requires an active internet connection to function.

\section{System Features}

\subsection{Upload Academic Resource}

\textbf{Description:} A student uploads academic materials such as lecture notes, project reports, or question papers to a subject repository for sharing.

\textbf{Priority:} High

\textbf{Main Flow:}
\begin{itemize}
    \item Student accesses the ``Contribute'' section.
    \item Selects category and uploads file (PDF, DOCX, PPT).
    \item Adds required metadata (Subject, Topic, Tags).
    \item System validates and publishes the resource.
\end{itemize}

\textbf{Key Requirements:}
\begin{itemize}
    \item Student must be logged in with Contributor permissions.
    \item Maximum file size is 50MB.
    \item Resource is stored and indexed in the database.
\end{itemize}


\subsection{Search and Download Resource}

\textbf{Description:} Students search academic materials using keywords or filters and download them for study.

\textbf{Priority:} Medium

\textbf{Main Flow:}
\begin{itemize}
    \item Student enters keywords or applies filters.
    \item System displays ranked results.
    \item Student downloads selected resource.
\end{itemize}

\textbf{Key Requirements:}
\begin{itemize}
    \item Student must be logged in.
    \item Download count is updated.
\end{itemize}


\subsection{View Tech-Trend Feed}

\textbf{Description:} Students view real-time technological updates related to their academic domains.

\textbf{Priority:} Medium

\textbf{Main Flow:}
\begin{itemize}
    \item Student accesses the Tech-Trend section.
    \item System fetches and displays curated updates.
\end{itemize}

\textbf{Requirement:}
\begin{itemize}
    \item Active internet connection is required.
\end{itemize}


\subsection{Rate and Review Resource}

\textbf{Description:} Students provide ratings and optional reviews for shared resources.

\textbf{Priority:} Low

\textbf{Main Flow:}
\begin{itemize}
    \item Student selects rating (1--5 stars).
    \item Optional comment is submitted.
    \item System updates the average rating.
\end{itemize}

\textbf{Requirement:}
\begin{itemize}
    \item Student must have accessed the resource.
\end{itemize}


\subsection{User Management}

\textbf{Description:} Nexus Notes administrators manage user roles and permissions.

\textbf{Priority:} High

\textbf{Main Flow:}
\begin{itemize}
    \item Administrator accesses the User Management section.
    \item Updates user roles or permissions.
    \item System saves changes.
\end{itemize}

\textbf{Requirement:}
\begin{itemize}
    \item Administrator must be logged in with appropriate privileges.
\end{itemize}


\subsection{Content Moderation}

\textbf{Description:} Nexus Notes administrators review and manage uploaded content.

\textbf{Priority:} High

\textbf{Main Flow:}
\begin{itemize}
    \item Administrator reviews flagged resources.
    \item Deletes or approves content.
    \item System updates content visibility.
\end{itemize}

\textbf{Requirement:}
\begin{itemize}
    \item Content must be flagged or under review.
\end{itemize}

\section{Other Nonfunctional Requirements}

\subsection{Performance Requirements}
The Nexus Notes platform shall operate efficiently under normal and peak usage conditions.

\begin{itemize}
    \item The system shall support concurrent access by multiple users without noticeable performance degradation.
    \item Search results shall be displayed within 2--3 seconds under normal load.
    \item File uploads up to 50MB shall be processed and validated efficiently.
    \item The platform shall be scalable to accommodate increasing numbers of users and resources.
    \item System availability shall remain high except during scheduled maintenance.
\end{itemize}

\subsection{Security and Privacy Constraints}
The system shall ensure protection of user data and academic resources.

\begin{itemize}
    \item User authentication shall be required to access system features.
    \item Role-based access control shall restrict actions based on user type (Student, Faculty, Administrator).
    \item Passwords and sensitive data shall be securely stored.
    \item Secure communication protocols shall be used for data transmission.
    \item The system shall comply with institutional data protection policies.
\end{itemize}

\subsection{Software Quality Attributes}
Nexus Notes shall follow standard software development practices to ensure reliability and maintainability.

\begin{itemize}
    \item The system shall provide a user-friendly and intuitive interface.
    \item The platform shall support responsive design for different devices.
    \item Periodic backups shall be maintained to prevent data loss.
\end{itemize}
\section{Other Requirements}

\subsection{Applicable Standards}
Database updates and backups ensure academic resources are securely managed.
The system provides a mechanism for continuous updates of digital content.

\end{document}